% Condensed from lit_review_condensed.md; ~3/4 column.
\paragraph{From passive to active spatial reasoning.}
Static-scene benchmarks \citep{weston2015babi,shi2022stepgame,
mirzaee2021spartqa,chen2024spatialvlm,ma20243dsr,yang2025mmsi} evaluate
disembodied spatial reasoning. A second wave introduced
\emph{cognitive-map prediction} as an auxiliary objective:
VSI-Bench \citep{yang2025vsi} and MindCube \citep{yin2025mindcube}
showed that explicit map prediction improves video and multi-view QA.
\tos \citep{zhang2026theoryofspace} departs from both lines: the agent
itself decides what to observe next, externalises its cognitive map at
every step (probing), and is tested under a \emph{false-belief}
perturbation borrowed from developmental psychology
\citep{wimmer1983false}. \tos provides the only published probing
benchmark that simultaneously evaluates Stability, Self-tracking,
Local$\leftrightarrow$Global consistency, and Belief Inertia under a
single representation.

\paragraph{Cognitive maps in neuroscience and bio-inspired AI.}
Place cells \citep{okeefe1971}, grid cells \citep{hafting2005}, and
their hierarchical structure \citep{stensola2012} encode space at the
region level rather than as a flat coordinate list. Banino~et~al.\
\citep{banino2018grid} showed that grid-like codes \emph{spontaneously
emerge} in deep RL agents; MoHA (Motivational Cognitive Maps,
\citet{mohacog2025}) modulates the spatial map by motivational
signals; RoboMemory \citep{robomemory2024} integrates spatial,
temporal, and episodic memory in a multi-memory architecture. None of
these works yet expose their region-level structure to a probing
benchmark of the kind \tos provides.

\paragraph{Territoriality, familiarity, ownership.}
Behavioural ecology defines \emph{territory} as the spatial unit at
which familiarity, defence, and recovery are organised
\citep{burt1943,powell2012homerange}. AI work has not directly
borrowed this framing: count-based exploration bonuses
\citep{burda2018rnd,pathak2017icm} treat visitation as a reward
signal, not a structural property of regions; Voronoi-based coverage
\citep{cortes2004coverage} and learned multi-robot patrolling
\citep{magec2024,patrolppo2025} pre-assign or emergently allocate
regions but never let the partition itself emerge from familiarity.
The \emph{dual-source territoriality} framing this project realises
---physical (room walls) plus familiarity (visit-weighted)---has no
direct precedent in the foundation-model literature.

\paragraph{World models and long-horizon spatial tasks.}
DreamerV3 \citep{hafner2023dreamerv3} and the Memory Maze benchmark
\citep{pasukonis2023memorymaze} show that latent world models
substantially outperform model-free methods on long-horizon memory.
Recall-to-Imagine \citep{r2i2024} adds episodic retrieval and reaches
superhuman Memory-Maze performance. These methods, however, treat the
latent space as monolithic; territorial factorisation of the latent
remains an explicit open direction. VAGEN \citep{wang2025vagen}---the
multi-turn VLM-RL framework on which the \tos environment is built---
provides the only ready infrastructure to actually \emph{train} a
world model on \tos. Training is out of scope for this 3-day course
project, but the post-hoc \tb representation we evaluate is shaped to
be a drop-in substrate for that follow-up.
